% !TEX root = ../master.tex

\section{\texttt{RecursiveSet}: Interne Architektur und Implementierung}
\label{sec:recursive-set-intern}

Während Kapitel 3 die fachliche API und die konzeptionelle Notwendigkeit einer wertsemantischen Mengenstruktur behandelte, fokussiert dieses Kapitel die zugrunde liegende technische Architektur. Da das native JavaScript-\mintinline{typescript}{Set} keine strukturelle Gleichheit für komplexe Objekte unterstützt und sich nicht durch einfache Wrapper modifizieren lässt, erfordert \mintinline{typescript}{RecursiveSet} eine vollständige Eigenimplementierung als Hash-Tabelle. 

Die zentrale Herausforderung besteht darin, strukturelle Gleichheitsprüfungen (\(O(1)\) amortisiert), effiziente Speicherverwaltung und Unveränderlichkeitsgarantien in einer hochperformanten, Garbage-Collection-freundlichen Architektur zu vereinen.

\subsection{Hash-Dispatcher und Gleichheitsprüfung}
\label{subsec:rs-intern-hashing}

Die technische Basis für den amortisierten \(O(1)\)-Zugriff bildet die strikte Trennung von Hash-Berechnung (zur Reduktion des Suchraums) und der tatsächlichen strukturellen Gleichheitsprüfung.

Die Funktion \mintinline{typescript}{getHashCode} agiert als Dispatcher, der für jeden nach dem \mintinline{typescript}{Value}-Vertrag (siehe Abschnitt \ref{subsec:rs-api-wertetypen}) zulässigen Typ eine spezialisierte Hash-Strategie wählt:

\begin{listing}[H]
	\caption{Optimierte Hash-Berechnung für numerische Typen}
	\label{lst:rs-gethashcode-number}
	\begin{minted}{typescript}
		if (typeof val === 'number') {
			// Fast Path: 32-Bit Integers (Smis)
			if ((val | 0) === val) {
				let h = val | 0;
				h = Math.imul(h ^ (h >>> 16), 0x85ebca6b);
				h = Math.imul(h ^ (h >>> 13), 0xc2b2ae35);
				return (h ^ (h >>> 16)) >>> 0;
			}
			// Slow Path: IEEE 754 Doubles (via ArrayBuffer)
			_f64[0] = val;
			return (Math.imul(_i32[0], 0xcc9e2d51) ^ Math.imul(_i32[1], 0x1b873593)) >>> 0;
		}
	\end{minted}
\end{listing}

Um Kollisionen bei sequenziellen Ganzzahlen zu vermeiden, wird der Wert durch Bit-Mixing (ähnlich zu MurmurHash3) gestreut. Fließkommazahlen werden über einen \mintinline{typescript}{ArrayBuffer} in ihre binäre IEEE-754-Repräsentation zerlegt, um FPU-Overhead zu vermeiden. Strings verwenden eine Jenkins-One-at-a-Time-Variante, während bei strukturellen Objekten direkt auf die vertraglich zugesicherte Eigenschaft \mintinline{typescript}{val.hashCode} delegiert wird.

Die finale Gleichheitsprüfung erfolgt in \mintinline{typescript}{areEqual}. Nach einem \mintinline{typescript}{===} Kurzschluss (Fast Path) greift die Methode für Zahlen auf \mintinline{typescript}{Object.is} zurück, um Edge-Cases wie \mintinline{typescript}{-0} und \mintinline{typescript}{NaN} sprachkonform zu isolieren, und delegiert bei komplexen Typen an \mintinline{typescript}{val.equals()}.


\subsection{Speicherlayout: Structure of Arrays (SoA)}
\label{subsec:rs-intern-speicherlayout}

Eine naive Implementierung einer Hash-Map würde Bucket-Objekte der Form \mintinline{typescript}{{ key, value, hash, next }} allozieren. Dies führt in der V8-Engine jedoch zu Speicherfragmentierung und hohen Latenzen durch die Garbage Collection. \mintinline{typescript}{RecursiveMap} (als architektonische Basis für \mintinline{typescript}{RecursiveSet}) nutzt stattdessen das \textit{Structure of Arrays} (SoA) Pattern.

\begin{listing}[H]
	\caption{Speicherarchitektur von \texttt{RecursiveMap}}
	\label{lst:recursivemap-storage-layout}
	\begin{minted}{typescript}
		class RecursiveMap<K extends Value, V extends Value> {
			// Dichte Arrays (Nutzdaten)
			protected readonly _keys: K[] = [];
			protected readonly _values: V[] = [];
			protected readonly _hashes: number[] = [];
			
			// Indextabelle (Suchstruktur für Open Addressing)
			protected readonly _indices: Uint32Array;
		}
	\end{minted}
\end{listing}

Die Nutzdaten wachsen dicht (ohne Lücken) in den Arrays \texttt{\_keys}, \texttt{\_values} und \texttt{\_hashes}. Die eigentliche Hash-Tabelle ist das \mintinline{typescript}{Uint32Array} namens \texttt{\_indices}. Es enthält keine Objekte, sondern \textbf{1-basierte Zeiger} auf die dichten Arrays. Ein Wert von \texttt{0} markiert einen leeren Slot. Ein Wert \(n > 0\) verweist auf den Index \(n - 1\) in den dichten Arrays.


\subsection{Kollisionsauflösung via Open Addressing}
\label{subsec:rs-intern-open-addressing}

Das Nachschlagen und Einfügen erfolgt über \textit{Open Addressing} mit \textit{Linear Probing}. Kollidieren Hash-Werte auf demselben Bucket-Slot, wird linear der nächste freie Slot gesucht. 

Der Suchalgorithmus profitiert maßgeblich vom SoA-Layout: Bevor die teure strukturelle Methode \mintinline{typescript}{areEqual} aufgerufen wird, vergleicht die Schleife den vorberechneten Hash aus \texttt{\_hashes}.

\begin{listing}[H]
	\caption{Lookup-Algorithmus mit linearem Sondieren}
	\label{lst:recursivemap-find-schema}
	\begin{minted}{typescript}
		get(key: K): V | undefined {
			const h = getHashCode(key);
			let idx = h & this._mask; // Start-Bucket berechnen
			
			while (true) {
				const entry = this._indices[idx];
				if (entry === 0) return undefined; // Slot leer -> Nicht gefunden
				
				const ptr = entry - 1;
				// 1. Hash vergleichen (schnell), 2. equals vergleichen (teuer)
				if (this._hashes[ptr] === h && areEqual(this._keys[ptr], key)) {
					return this._values[ptr];
				}
				idx = (idx + 1) & this._mask; // Nächster Slot (Linear Probing)
			}
		}
	\end{minted}
\end{listing}


\subsection{Löschvorgang: Backshift Deletion und Dense-Swap}
\label{subsec:rs-intern-backshift}

Das Entfernen von Elementen aus Hash-Tabellen mit Open Addressing ist komplex, da das einfache Leeren eines Slots die Suchketten (Probe Chains) nachfolgender Elemente durchtrennen würde. Die übliche Lösung – das Setzen von "Tombstones" (als gelöscht markierte Slots) – führt langfristig zu verschmutzten Tabellen und Performanceverlust. 

Die Bibliothek verwendet stattdessen \textbf{Backshift Deletion} (auch Robin Hood Deletion genannt). Die Lücke im \texttt{\_indices}-Array wird durch nachfolgende Elemente aufgefüllt, sofern deren idealer Hash-Bucket weiter vorne liegt.

\begin{listing}[H]
	\caption{Reparatur der Probe Chain (Backshifting)}
	\label{lst:rs-backshift}
	\begin{minted}{typescript}
		private removeIndex(holeIdx: number): void {
			let i = (holeIdx + 1) & this._mask;
			while (this._indices[i] !== 0) {
				const entry = this._indices[i];
				const h = this._hashes[entry - 1]; 
				
				const idealIdx = h & this._mask;
				const distToHole = (holeIdx - idealIdx + this._bucketCount) & this._mask;
				const distToI = (i - idealIdx + this._bucketCount) & this._mask;
				
				// Element i ist näher an seinem idealen Slot, wenn es nach holeIdx rückt
				if (distToHole < distToI) {
					this._indices[holeIdx] = entry;
					holeIdx = i;
				}
				i = (i + 1) & this._mask;
			}
			this._indices[holeIdx] = 0; // Lücke ans Ende der Chain verschoben und geleert
		}
	\end{minted}
\end{listing}

Nachdem die Suchstruktur repariert ist, muss das Element aus den dichten Nutzdaten (\texttt{\_keys}, \texttt{\_values}, \texttt{\_hashes}) entfernt werden. Ein \mintinline{typescript}{Array.prototype.splice()} würde ein \(\mathcal{O}(N)\)-Verschieben aller nachfolgenden Elemente erzwingen. Stattdessen nutzt die Implementierung den \textbf{Dense-Swap}: Das zu löschende Element wird mit dem \textit{letzten} Element der dichten Arrays überschrieben. Anschließend wird der Eintrag des verschobenen letzten Elements in der \texttt{\_indices}-Tabelle auf die neue Position aktualisiert. Dies garantiert eine \(O(1)\)-Löschung ohne Array-Reallokation.


\subsection{Kapazitätsverwaltung und Hashing-Metadaten}
\label{subsec:rs-intern-resize-shrink}

Um die O(1)-Performance aufrechtzuerhalten, wird der Auslastungsgrad (Load Factor) strikt bei \(0.75\) (\(75\%\)) abgeriegelt. Wird dieser überschritten, verdoppelt \mintinline{typescript}{ensureCapacity} die Tabellengröße (wobei die Kapazität stets eine Potenz von 2 bleibt, um Modulo-Operationen durch bitweises UND \mintinline{typescript}{& this._mask} zu ersetzen).

Um Speicherlecks in stark fluktuierenden Datenstrukturen zu vermeiden, greift bei Deletionen der \textit{Cold Path} (\mintinline{typescript}{shrink()}): Fällt die Auslastung unter \(25\%\) bei einer Mindestgröße von \texttt{MIN\_BUCKETS = 16}, wird die Indextabelle halbiert und re-hasht.


\subsection{\texttt{RecursiveSet}: Architektonische Spezialisierung}
\label{subsec:rs-intern-recursiveset}

\mintinline{typescript}{RecursiveSet} verwendet dieselbe Speicher- und Sondierungsarchitektur wie die Map, eliminiert jedoch das \texttt{\_keys}-Array, da die gespeicherten Werte zeitgleich ihre eigenen Suchschlüssel sind. 

Die Besonderheit des \mintinline{typescript}{RecursiveSet} liegt in seiner eigenen \mintinline{typescript}{hashCode}-Generierung. Da Mengen in der Theorie ungeordnet sind, muss ihr Hash unabhängig von der Einfügereihenfolge sein. Dies wird durch die kumulative bitweise Exklusiv-Oder-Verknüpfung (\mintinline{typescript}{^=}) der Element-Hashes im Feld \mintinline{typescript}{this._xorHash} erreicht.

\paragraph{Immutabilität und der Freeze-Zyklus}
\label{subsec:rs-intern-freeze}

Wie in Abschnitt \ref{subsec:rs-api-freeze-frozenset} erläutert, verbindet das Set die Phasen der Mutabilität und Immutabilität. Technisch wird dies durch den Getter \mintinline{typescript}{get hashCode()} erzwungen. Sobald eine Menge (z. B. als Element einer anderen Menge) gehasht wird, setzt sie intern \mintinline{typescript}{this._isFrozen = true}. 

Zudem greift ein kaskadierender Freeze: Der Getter iteriert über alle gespeicherten komplexen Objekte und ruft deren \mintinline{typescript}{hashCode} auf. Dies garantiert, dass auch tief verschachtelte Mengenstrukturen (\mintinline{typescript}{RecursiveSet<RecursiveSet<T>>}) iterativ bis in die unterste Ebene eingefroren und versiegelt werden.


\subsection{Performance-Optimierung funktionaler Methoden}
\label{subsec:rs-intern-functional-methods}

Die in Kapitel 3 eingeführten Transformationen (\mintinline{typescript}{map}, \mintinline{typescript}{filterMap}, \mintinline{typescript}{flatMap}) basieren technisch nicht auf Iteratoren oder generischen JavaScript-Array-Methoden. Um den Overhead von Generator-Funktionen und Callbacks zu minimieren, iterieren diese Methoden imperativ direkt über das dichte \texttt{\_values}-Array.

Zudem wird bei \mintinline{typescript}{map} und \mintinline{typescript}{filterMap} eine optimistische Prä-Allokation durchgeführt: \mintinline{typescript}{result.ensureCapacity(this.size)}. Unter der Annahme, dass die transformierte Menge ähnlich groß wird wie die Ursprungsmenge, werden teure \mintinline{typescript}{resize()}-Zyklen während der Iteration vollständig eliminiert.


\subsection{Fazit der internen Architektur}
\label{subsec:rs-intern-fazit}

Die interne Umsetzung demonstriert, dass Wertsemantik in TypeScript nicht durch bloße Typ-Deklarationen, sondern nur durch eine tiefgreifende Modifikation der Speicher- und Zugriffsstrukturen erreicht werden kann. 

Durch das Structure of Arrays-Paradigma (SoA), die Trennung von Hash-Wegweiser und Gleichheitsprüfung, sowie Algorithmen wie Backshift-Deletion und Dense-Swapping wird \mintinline{typescript}{RecursiveSet} zu einer hochperformanten Infrastruktur. Sie bietet die formale Strenge unveränderlicher mathematischer Mengen und liefert gleichzeitig die notwendige Ausführungsgeschwindigkeit, um komplexe deterministische Algorithmen der theoretischen Informatik in Laufzeitumgebungen wie V8 abzubilden.